\documentclass[journal]{IEEEtran}  % IEEE期刊模板
\usepackage{amsmath,amsfonts}      % amsmath提供多行公式对齐、矩阵环境等数学排版功能；amsfonts扩展数学符号库‌
\usepackage{algorithmic}          % 实现算法伪代码环境，支持IF/WHILE等结构化语句排版‌
\usepackage{algorithm}            % 配合algorithmic使用，提供浮动算法环境及标题控制‌
\usepackage{graphicx}              % 图像插入核心工具，支持主流图片格式‌
\usepackage{cite}                  % 智能处理参考文献引用排序与压缩格式‌
\hyphenation{op-tical net-works semi-conduc-tor IEEE-Xplore} %是LaTeX 中用于‌自定义单词断词规则‌的命令
\graphicspath{{figures/}}  % 图片的路径
\usepackage[unicode]{hyperref} %提供跳转链接

\begin{document}

\title{Secret}

\author{Enjun Xia
        % <-this % stops a space
\thanks{Manuscript received xxxx, 2023; revised xxxx, 2023. This work was supported by . }% <-this % stops a space
\thanks{The authors are with the School of (emails: ).}
}

% The paper headers
\markboth{IEEE ,~Vol.~xx, No.~xx, xxxx,~2023}%
{Xia \MakeLowercase{\textit{et al.}}: }

%\IEEEpubid{0000--0000/00\$00.00~\copyright~2021 IEEE}
% Remember, if you use this you must call \IEEEpubidadjcol in the second
% column for its text to clear the IEEEpubid mark.

\maketitle

\begin{abstract}

\end{abstract}

\begin{IEEEkeywords}

\end{IEEEkeywords}

\section{Introduction}
\IEEEPARstart{D}{ifferent} from



\bibliographystyle{IEEEtran}
\bibliography{myref}


\end{document}


